\begin{python0}
from solutions import *; clear()
\end{python0}
\sectionthree{Refactoring}

Note that the program prompts the user for 4 int values and displays an expression for it. We could also reorganize the program like this:

\begin{consolethree}[escapeinside=||]
void get_user_input(int & a, int & b, int & c, int & d)
{
    std::cin >> a >> b >> c >> d;
}

void display_expression(int a, int b, int c, int d)
{
    std::cout << a << " + " << b << " + " << c << " + " << d << '\n';
}

int main()
{
    int a, b, c, d;
    get_user_input(a, b, c, d);
    display_expression(a, b, c, d);
    return 0;
}
\end{consolethree}

Remember that the new program has only changed in terms of organization (with the addition of functions).

The above exercise where you modify the structure of a code without changing its functionality is called \EMPHASIZE{refactoring}.

Sometimes the changes are ad hoc. Without constant refactoring, software can decay into spaghetti code that no one can understand.

